\clearpage
\begingroup
\fancyhead[L]{\small POSTHUMOUS PAPERS AND CORRESPONDENCE}
\renewcommand{\headrulewidth}{0pt}
\SourceLeaf{605}
\par\medskip
\Needspace{8\baselineskip}
\pdfbookmark[2]{[7]}{pallas-IX-fragment7}
\begin{center}[7.]\end{center}
\begingroup
\small
\par\noindent\hspace*{\parindent}[The quantities \(M,N,\varphi\) must be calculated by approximation from equations 3 in
note [5]. The roots of the
equation \(\rho\rho=0\) can also be determined in the following way.]
\par\noindent\hspace*{\parindent}Solution of the equation
\[
a+b\cos2u=\cos(u+D)
\]
\noindent[where \(u\) is written for \(E'\) and the relations of the constants \(a,b,D\)
to those in equation 1 of note [5] are readily apparent].
\par\noindent\hspace*{\parindent}Set
\[
u=x+y,
\]
\noindent so that
\[
\begin{aligned}
a+b\cos2x\cos2y-\cos(D+x)\cos y&\\
{}+\sin y(\sin(D+x)-2b\sin2x\cos y)&=0.
\end{aligned}
\]
\par\noindent\hspace*{\parindent}If, then, as we may, we set
\[
\text{1)}\qquad 2b\sin2x\cos y=\sin(D+x),
\]
\noindent we obtain
\[
\text{2)}\qquad b\cos2x\cos2y-\cos(D+x)\cos y+a=0.
\]
\par\noindent\hspace*{\parindent}Substituting the value of \(\cos y\) from 1) into 2) gives
\[
4b\sin2x^2\mathord{.}(a-b\cos2x)=\cos2D-\cos2x,
\]
\noindent so, setting \(\cos2x=t\), we have
\[
\text{3)}\qquad 4b(1-t)(1+t)(a-bt)=\cos2D-t
\]
\noindent or, on expanding,
\[
\text{4)}\qquad 4bbt^3-4abtt+(1-4bb)t+4ab-\cos2D=0.
\]
\par\noindent\hspace*{\parindent}Form 3) is the most convenient for an indirect solution.
\par\noindent\hspace*{\parindent}Once \(x\) is found, \(y\) can be determined either from 1) by
the formula
\[
\cos y=\frac{\sin(D+x)}{2b\sin2x}
\]
\SourceLeaf{606}
\noindent or from 2) by the formula
\[
\cos y=\frac{(-a+b\cos2x)\sin2x}{\sin(D-x)}
\]
\noindent For \(x\), one can be certain of always obtaining a real value.
If a value greater than 1 is found for \(\cos y\), the imaginary
value of \(y\) is
\[
y=i\log\mathrm{hyp}\{\cos y+\BareRoot{(\cos y^2-1)}\}.
\]
\noindent If \(u\) denotes an eccentric anomaly and \(e\) the eccentricity,
the corresponding mean anomaly is
\[
\begin{aligned}
x-e\sin x\cos y+i\{&\log\mathrm{hyp}(\cos y+\BareRoot{(\cos y^2-1)})\\
&-e\cos x\BareRoot{(\cos y^2-1)}\}.
\end{aligned}
\]
\par\noindent\hspace*{\parindent}If
\[
a\cos(nM+A),\qquad a'\cos((n+1)M+A')
\]
\noindent are two consecutive terms of a series representing a function that becomes
infinite when \(a+b\cos2u=\cos(u+D)\), then
\[
\begin{gathered}
\mathrm{Limes}\,(A-A')=x-e\sin x\cos y\\
\mathrm{Limes}\,\frac a{a'}=\frac{\cos y+\BareRoot{(\cos y^2-1)}}{e\cos x\BareRoot{(\cos y^2-1)}},
\end{gathered}
\]
\noindent where \(c\) denotes the base of the hyperbolic logarithms. To make
the calculation easier, we can also set \(\cos y=\frac1{\sin\zeta}\), which gives
\[
\mathrm{Limes}\,\frac a{a'}=\frac{\mathrm{cotang}\,\frac12\zeta}{c^{e\cos x\,\mathrm{cotang}\,\zeta}}
\]
\noindent for the limit.
\par\noindent\hspace*{\parindent}[Thus, as above, for the factor
\(\lambda=\left(\frac1{\sin\zeta'}-\cos(M'-\varphi')\right)^{\frac32}\) we take the values
corresponding to the root of the equation \(\lambda=0\)
\[
\varphi'=x-\frac{e\sin x}{\sin\zeta}=\mathrm{Limes}\,(A-A')
\]
\[
\begin{gathered}
\frac1{\sin\zeta'}=\cos\left(i\log\mathrm{hyp}\,\mathrm{Limes}\,\frac a{a'}\right)\\
\text{or}\qquad \mathrm{cotang}\,\tfrac12\zeta'=\mathrm{Limes}\,\frac a{a'}\mathpunct{\cdot}\Biggr]
\end{gathered}
\]
\par\medskip\begin{center}\rule{4em}{.4pt}\end{center}
\endgroup

\clearpage
\endgroup
